\subsection{Task 8: Tỷ lệ lấp đầy theo tháng và chính sách số đêm tối thiểu}

\subsubsection{Thiết kế Idiom}

\begin{table}[H]
\centering
\renewcommand{\arraystretch}{1.1}
\begin{tabularx}{\textwidth}{|l|l|X|}
\hline
\textbf{Idiom} & \multicolumn{2}{l|}{\textbf{Heatmap (Matrix Chart)}} \\ \hline
\textbf{What}
& \multicolumn{2}{X|}{Month: Ordinal (tháng 1--12, cyclic)} \\ \cline{2-3}
& \multicolumn{2}{X|}{minimum\_nights\_group (derived): Categorical (Short $\leq$3 / Medium 4--7 / Long $>$7)} \\ \cline{2-3}
& \multicolumn{2}{X|}{occupancy\_rate\_pct (derived): Quantitative} \\ \hline
\textbf{How}
& \textbf{Encode} &
\textbf{Mark:} Point (Square/Area lấp đầy ô lưới ma trận) \newline
\textbf{Channel:} \newline
\hspace*{1em} Pos X: Month (Ordinal -- 12 cột) \newline
\hspace*{1em} Pos Y: Minimum Nights Group (Categorical -- Short/Medium/Long, top$\rightarrow$bottom) \newline
\hspace*{1em} Color (Luminance/Sequential): Occupancy Rate (Quantitative -- đậm=cao, nhạt=thấp) \\ \cline{2-3}
& \textbf{Manipulate} &
Hover để xem month, Minimum Nights Group, occupancy\_rate\_pct (\%). \\ \cline{2-3}
& \textbf{Reduce} &
Binning minimum\_nights $\rightarrow$ 3 nhóm (giảm số chiều từ hàng trăm giá trị $\rightarrow$ 3 nhóm). \newline
Sort thủ công: Short $\rightarrow$ Medium $\rightarrow$ Long (thứ tự logic nghiệp vụ, không alphabetical). \\ \hline
\textbf{Why} & \multicolumn{2}{X|}{
\textbf{discover $\rightarrow$ compare} \newline
Phát hiện pattern 2 chiều đồng thời (month $\times$ minimum\_nights\_group). Heatmap theo nguyên lý `2 keys $\rightarrow$ heatmap' --- encode 3 biến trong cấu trúc ma trận. Color Luminance/Sequential đúng cho quantitative vì kích thước ô đồng đều loại bỏ bias từ Area.
} \\ \hline
\textbf{Scale} & \multicolumn{2}{X|}{
Key: 12 (tháng) $\times$ 3 (nhóm) = 36 ô. \newline
Bin: 3 (Short $\leq$3, Medium 4--7, Long $>$7 đêm).
} \\ \hline
\end{tabularx}
\end{table}

\begin{figure}[H]
    \centering
    \safeincludegraphics[width=0.9\linewidth]{charts/domain_task8_chart2.png}
    \caption{Hình 8.2. Tỷ lệ lấp đầy theo tháng và nhóm số đêm tối thiểu}
    \label{fig:domain-task8-chart2}
\end{figure}

\subsubsection{Đánh giá biểu đồ}

\textbf{1. Nguyên lý biểu đạt (Expressiveness):}

\begin{itemize}
    \item Thuộc tính: Tháng / Month (O), Nhóm đêm tối thiểu / Minimum Nights Group (C -- derived: Short $\leq$3 / Medium 4--7 / Long $>$7), Tỷ lệ lấp đầy / Occupancy Rate Pct (Q -- Color Luminance).
    \item Channel:
    \begin{itemize}
        \item PosX (Month -- O): trục thời gian theo thứ tự $\rightarrow$ đúng.
        \item PosY (Minimum Nights Group -- C): phân loại chính sách $\rightarrow$ đúng. Spatial region cho categorical.
        \item Color (Luminance/Sequential): thể hiện độ lớn occupancy (Q) $\rightarrow$ phù hợp. Sequential colormap (nhạt$\rightarrow$đậm) đúng cho quantitative attribute có hướng sequential.
    \end{itemize}
    \item Đánh giá mức độ quan trọng:
    \begin{itemize}
        \item Ưu tiên 1 -- Vị trí ngang (PosX --- chiều thời gian Ordinal): PosX mã hóa Month (Ordinal --- tháng 1 đến 12) dọc theo trục ngang của ma trận heatmap. Đây là kênh quan trọng nhất vì chiều thời gian là trục phân tích chính của task (seasonality). Theo Mackinlay, Position là kênh mạnh nhất cho dữ liệu thứ tự (Ordinal), đảm bảo thứ tự tháng được duy trì chính xác và người xem có thể theo dõi sự thay đổi của occupancy theo thời gian một cách tự nhiên từ trái sang phải --- phù hợp với convention đọc phương Tây và nhận thức thời gian tuyến tính.
        \item Ưu tiên 2 -- Vị trí dọc (PosY --- Spatial Region cho Categorical): PosY mã hóa Minimum Nights Group (Categorical --- Short/Medium/Long stay) theo hàng của ma trận. Đây là kênh quan trọng thứ hai vì nó xác định chiều so sánh thứ hai của biểu đồ: người xem so sánh giữa các chính sách minimum nights theo trục dọc. Theo Mackinlay, Spatial Region (phân vùng không gian) phù hợp cho dữ liệu danh mục --- mỗi hàng heatmap đóng vai trò như một spatial region riêng biệt, phân tách rõ ràng ba nhóm chính sách mà không cần dùng màu sắc (vốn đã được dùng cho kênh thứ ba).
        \item Ưu tiên 3 -- Độ sáng màu sắc (Color Luminance / Sequential): Color Luminance mã hóa Occupancy Rate (Quantitative) thông qua thang màu tuần tự (nhạt $\rightarrow$ đậm). Theo Mackinlay, Color Luminance không phải kênh tối ưu nhất cho dữ liệu định lượng về độ chính xác tuyệt đối (kém hơn Position hay Length), nhưng trong heatmap --- nơi cả PosX và PosY đã được dùng để mã hóa hai biến phân loại/thứ tự --- đây là kênh duy nhất khả dụng để mã hóa giá trị định lượng thứ ba trong không gian hai chiều. Color Luminance phù hợp để phát hiện pattern (nhận biết ô đậm/nhạt) và so sánh tương đối giữa các ô, phục vụ đúng mục tiêu discover của heatmap. Kích thước ô đồng đều (square) loại bỏ nhiễu từ kênh Area, giúp người xem chỉ tập trung vào màu sắc.
    \end{itemize}
    \item Kết luận: Heatmap là idiom phù hợp nhất khi cần phân tích 3 biến (2 keys + 1 value) trong cấu trúc ma trận.
\end{itemize}

\textbf{2. Nguyên lý hiệu quả (Effectiveness):}

\begin{itemize}
    \item Độ chính xác (Accuracy): Color (Luminance) --- accuracy thấp hơn Position, khó ước lượng chính xác \% chênh lệch. Tuy nhiên mục tiêu là phát hiện pattern (cao/thấp), không phải đọc giá trị chính xác --- trade-off chấp nhận được.
    \item Khả năng phân biệt (Discriminability): 36 ô vuông đều nhau. Sequential palette giúp phân biệt $\sim$5--7 cấp độ màu.
    \item Khả năng tách biệt (Separability): PosX, PosY và Color kết hợp tốt. Area đồng đều không tạo bias về kích thước.
\end{itemize}

\subsubsection{Phân tích biểu đồ}

\begin{itemize}
    \item Short minimum stay ($\leq$3 đêm) có occupancy cao và ổn định quanh năm --- chính sách linh hoạt thu hút nhiều khách nhất, đặc biệt trong mùa thấp điểm.
    \item Long minimum stay ($>$7 đêm) có occupancy thấp hơn trong phần lớn các tháng, nhưng tăng vào cuối năm --- phù hợp cho khách công tác dài hạn.
    \item Tháng 11--12 có màu đậm nhất ở hầu hết các nhóm --- xác nhận high season với nhu cầu cao bất kể chính sách minimum nights.
    \item Khuyến nghị cho host: Tháng 1--4 nên chuyển sang Short minimum stay; tháng 11--12 có thể giữ Medium/Long stay vì cầu cao tự nhiên bù đắp tính hạn chế của chính sách.
\end{itemize}
